\PassOptionsToPackage{quiet}{xeCJK}
%% === v1.4.4 借鉴自 数模陪跑独家 模板：教学开关 + 占位符宏 ===
%% 在文档任意位置可用 \Teach{...} 写红色填写提示，\TeachExample{...} 写红色示例。
%% \Placeholder{...} 用黑色加粗显示【】占位符，便于交稿前搜索替换。
%% 默认 \showteachingfalse：教学提示不显示（不影响现有 PDF 输出）。
%% 想看教学提示：把下面 \showteachingfalse 注释掉，把上面 \showteachingtrue 取消注释即可。
\newif\ifshowteaching
\showteachingfalse
% \showteachingtrue

\newcommand{\Placeholder}[1]{\textbf{#1}}

\newcommand{\Teach}[1]{%
  \ifshowteaching
    \par\begingroup
      \noindent\color{red!78!black}\bfseries
      【填写提示（提交前自动隐藏）：#1】\par
    \endgroup
  \fi
}

\newcommand{\TeachExample}[1]{%
  \ifshowteaching
    \par\begingroup
      \noindent\color{red!78!black}\bfseries
      【示例（仅供理解，提交前自动隐藏）：#1】\par
    \endgroup
  \fi
}
%% === v1.4.4 借鉴结束 ===


%% 2026 规范修订：
%%   - 论文版（纸质版）= withpreface → 生成 承诺书(第1页) + 编号专用页(第2页) + 摘要页(第3页,页码1) + 正文
%%   - 电子版 = withoutpreface → 跳过承诺书和编号专用页,第一页直接是摘要页
%% 编译两份: xelatex 论文.tex  → 论文.pdf (纸质版提交) ; xelatex 电子版.tex  → 电子版.pdf (电子版提交)
\documentclass[withpreface,bwprint]{format}
\usepackage{etoolbox}
\usepackage{ctex}
\BeforeBeginEnvironment{tabular}{\zihao{-5}}
\usepackage[framemethod=TikZ]{mdframed}
\usepackage{url}
\usepackage{array}
\usepackage{tabularx}
\usepackage{longtable}
\newcolumntype{C}{>{\centering\arraybackslash}X}
\newcolumntype{R}{>{\raggedleft\arraybackslash}X}
\newcolumntype{L}{>{\raggedright\arraybackslash}X}

\renewcommand{\textbf}[1]{{\song\bfseries #1}}

\usepackage{tikz}
\usetikzlibrary{arrows.meta}
\tikzset{
  box/.style={rectangle, draw, minimum width=3.2cm, minimum height=0.9cm, align=center},
  arrow/.style={thick, -{Stealth}}
}

\title{烟幕干扰弹的投放策略}

\begin{document}

%% \maketitle 触发 withpreface 路径：生成承诺书 → 编号专用页 → 标题页（页码复位为 1）
\maketitle

\begin{abstract}

%% ==================== 通用摘要模板 (v1.5.5 重写, 通用化) ====================
%%
%% 设计原则:
%% 1. 整段都是 \TeX 注释 (以 % 开头), check 15 跳过 % 行 → dryrun 不报 RED
%% 2. 留 1 行 INLINE 占位符 【】, check 15 抓【】→ 跑题用户填完才 GREEN
%% 3. 跑题用户复制本文件后, 删注释 + 填摘要正文, 删 【】占位符行
%%
%% 字数: 800-1000 字, 1 页 A4
%% 结构 (v1.5.3 板块自查 01 借鉴 BZD 7 维诊断):
%%   - 第 1 段 (2 句): 背景 + 核心问题 (不堆砌数据/算法)
%%   - 主体 N 段: 一问一段, 每段含 [任务 → 数据/假设 → 模型+数学类型 → 算法 → 量化结果 → 结论]
%%   - 最后段: 创新点 + 检验方式 + 优势 + 推广价值 + 局限 + 改进方向
%%   - 关键词: 4-6 个, 3 方向 (问题/模型/算法创新)
%% 写法: 第三人称, 不写图表/公式/参考文献编号, 数值必须来自正文不得虚构

%%\textbf{第一段 (引言 2 句):} 随着【你的研究背景, 1 句】, 【你的研究对象/核心问题, 1 句】。本文基于【数据概况, 1082 例】开展【研究目标, 4 问研究】。

%%\textbf{针对问题一 (按你实际题改):} 为【任务, 1 句】, 采用【模型名 + 数学类型, 如: 多元线性回归 (MLR) + 随机森林回归 (RFR)】。结果表明【关键特征 + 量化结果】。

%%\textbf{针对问题二 (按你实际题改):} 以【异常判据】为基准, 建立【模型】。推荐【最佳时点/参数】。

%%\textbf{针对问题三 (按你实际题改):} 按【严格按题面给的分界, 如: 题面 [20,28) [28,32) [32,36) [36,40) ≥40】分层, 采用【模型】。结果表明【关键结论】。

%%\textbf{针对问题四 (按你实际题改):} 对【数据, 严格按题面指定 sheet】, 构建【模型】。推荐【决策规则 + 比例】。

%%\textbf{创新点:} (1) 【模型层创新】; (2) 【方法层创新】。经【检验方式】, 模型有较好的稳健性; 局限: 【1 句】。

%%\keywords{【关键词 1, 问题方向】;【关键词 2, 模型方向】;【关键词 3, 算法方向】;【关键词 4, 创新点】;【关键词 5, 应用方向】}

\textbf{【在这里粘贴你的 800-1000 字摘要正文, 严格按上面 6 段结构 (引言 + 4 问 + 创新) + 关键词。删掉这一行, 替换为你的实际摘要。本行留【】= dryrun check 15 RED, 必须替换。】}

\label{abstract:end}

\end{abstract}

\thispagestyle{plain}

%% 2026 规范第三条：第 3 页为摘要专用页，第四条：第 4 页起正文
%% 摘要无论是否占满 1 页都要独占这一页，正文另起一页
\newpage

%% 2026 规范第四条：论文不要目录，第四页起直接是正文，正文 ≤ 30 页
%% （2024 旧版曾有 \tableofcontents + \newpage，2025 修订稿已删除）

%% ==================== 问题重述 模板 (v1.5.6, BZD 数模社 2026 规范) ====================
%%
%% 借鉴 BZD 数模社 第一章 模板: 1.1 研究背景 + 1.2 问题回顾 + 1.3 研究综述
%%
%% 写作原则 (BZD 4 色教学法):
%% - 黑色 = 文章主体, 直接套用
%% - 蓝色 = 操作步骤
%% - 红色 = 解释说明
%% - 绿色 = 常见误区
%% - 黄色 = AI 提示词
%% - 灰色 = 自查清单
%%
%% 1.1 研究背景: 1-3 段, 研究对象 + 现实情境 + 主要矛盾 + 研究价值
%% 1.2 问题回顾: 按题号重述每问, 保留关键数据/单位/约束
%% 1.3 研究综述: ~500 字, 3-4 段, 现有方法 + 不足 + 本文切入点
%%
%% 跑题用户必须:
%% 1. 删掉所有 % 注释 (整段 L1-L29)
%% 2. 替换【】占位符为自己的实际内容
%% 3. 删 下面的 \AbstractSlot{...} 行 (check 15 抓这个 inline 【】)
%% 4. 跑题后用 LaTeX 编译一次验证 (xelatex 论文.tex × 2)

\section{问题重述}

\subsection{研究背景}

\textbf{【研究对象】}是指【简要定义, 1 句】, 在【应用领域, 1 句】中具有【主要作用, 1 句】。随着【社会/经济/技术/政策背景, 1 句】的发展, \textbf{【研究对象】}呈现出【主要趋势, 1 句】。

与此同时, 其发展仍受到【主要矛盾或现实问题, 1 句】的影响, 因此有必要对【核心研究问题, 1 句】进行系统研究。

\subsection{问题回顾}

根据题目提供的数据、条件和任务要求, 需要解决以下问题:

（1） 分析【问题一研究对象, 1 句】, 确定【需要识别/评价/计算的内容, 1 句】。

（2） 根据【数据条件, 1 句】, 建立模型完成【问题二的预测/评价/优化任务, 1 句】。

（3） 考虑【新增条件, 1 句】, 研究【问题三的主要任务, 1 句】。

（4） 结合【指定情境或可用的前述结果, 1 句】, 完成【问题四的任务, 1 句】。

若题目还有其他问题, 按照相同方式继续准确表述。

\subsection{研究综述}

现有研究主要从【方向一】、【方向二】和【方向三】分析该问题。在【方向一】方面, 常用【模型或方法】研究【具体任务】; 在【方向二】方面, 主要采用【模型或方法】分析【具体问题】; 在【方向三】方面, 相关研究通常通过【模型或方法】评价【研究对象】。

现有方法为研究该问题提供了基础, 但在【现有不足 1】、【现有不足 2】和【现有不足 3】方面仍具有一定局限。基于题目条件和研究任务, 有必要进一步从【本文建模角度】开展分析。

% 跑题后必须删下面这一行 (check 15 占位符自检)
\textbf{【这里粘贴从 1.1 到 1.3 拼起来的不超过 2 页内容, 删掉这一行】}

%% ==================== 问题分析 模板 (v1.5.6, BZD 数模社 2026 规范) ====================
%%
%% 借鉴 BZD 数模社 第二章 模板: 5 段连贯 + 整体求解框架图
%% 注意: BZD 建议不要再机械分 "2.1 问题一分析、2.2 问题二分析" 等孤立小节,
%%       而是用 5 段自然段体现问题联动.
%%
%% 第 1 段 (总体分析): 概括任务 + 问题关系 (递进/共享数据/参数传递)
%% 第 2 段 (数据分析): 数据来源 + 质量问题 + 预处理方法 + 选择理由
%% 第 3 段 (前序问题建模思路): 第 1 问的数学类型 + 模型选择 + 备选方案
%% 第 4 段 (后续问题承接): 后问的预测/优化/评价/情景分析
%% 第 5 段 (检验与输出): 检验方法 + 最终输出形式
%% 末尾: 整体求解框架图 (用 TikZ 或 Mermaid 跨问题流程图)
%%
%% 跑题用户必须:
%% 1. 删掉所有 % 注释 (整段 L1-L25)
%% 2. 替换【】占位符为自己的实际内容
%% 3. 删 下面的 \ProblemAnalysisSlot{...} 行 (check 15 抓这个 inline 【】)

\section{问题分析}

\textbf{【第 1 段: 总体分析, 1 段】} 题目围绕【研究对象, 1 句】提出了【概括全部任务, 1 句】。从数学建模角度看, 各问题分别涉及【评价/预测/分类/优化/决策等类型】, 并通过【数据/指标/参数/中间结果】形成较强的逻辑联系。前序问题主要用于完成【基础任务】, 所得【结果类型】将为后续【任务】提供输入; 后续问题则进一步考虑【新增条件】, 最终形成由【起点任务】到【终点任务】的完整分析过程。

\textbf{【第 2 段: 数据分析, 1 段】} 题目提供的数据主要包括【数据类型, 1 句】, 同时需要补充收集【外部数据, 1 句】。考虑到数据可能存在【缺失/异常/量纲不同/统计口径不一致/指标相关性较强】等问题, 需要先通过【数据处理方法, 1 句】完成数据清洗与口径统一, 再采用【标准化/相关性分析/主成分分析/特征筛选方法】提取有效信息, 为后续建模提供可靠数据基础。

\textbf{【第 3 段: 前序问题建模思路, 1 段】} 对于前序问题, 需要解决【任务概括, 1 句】, 其本质属于【数学问题类型, 如 评价/预测/分类/优化】。根据【数据特点和题目条件】, 可采用【主模型, 如 多元线性回归/随机森林/灰色预测/层次分析/差分进化等】刻画【变量之间的关系】, 并将得到的【权重/参数/关键指标/分类结果】传递至后续问题。为降低单一模型的局限, 可以选择【备选方法 1】和【备选方法 2】, 并依据【评价标准, 如 MAE/RMSE/AUC/约束满足率】进行比较。

\textbf{【第 4 段: 后续问题承接与扩展, 1 段】} 在前述分析的基础上, 后续问题需要进一步完成【预测/优化/影响评价/情景分析】。针对【预测任务】, 可选取【候选方法】分别刻画线性、非线性或时序变化规律, 并依据【误差指标】比较模型表现或构建组合预测; 针对【优化或决策任务】, 可将前述预测值或评价结果作为模型参数, 以【优化目标】建立【规划模型】, 同时设置【主要约束】; 对于【外部政策/环境变化/不确定性】, 可通过参数调整和情景设置分析其影响。

\textbf{【第 5 段: 检验与输出, 1 段】} 为检验模型的合理性, 可采用【误差分析/交叉验证/对比实验/灵敏度分析/稳健性检验】评价模型表现。最终输出【题目要求的预测值/评价结果/优化方案/建议】, 总体分析与求解过程如图 \ref{fig:overall_route} 所示。

\begin{figure}[htbp]
  \centering
  \begin{tikzpicture}[
    node distance=1.0cm and 1.2cm,
    box/.style={rectangle, draw, minimum width=2.8cm, minimum height=0.8cm, align=center},
    arrow/.style={->, >=stealth, thick}
  ]
    \node[box] (data) {数据预处理};
    \node[box, below=of data] (q1) {问题1: 因素分析};
    \node[box, below=of q1] (q2) {问题2: 时点优化};
    \node[box, below=of q2] (q3) {问题3: 分层优化};
    \node[box, below=of q3] (q4) {问题4: 集成决策};
    \node[box, below=of q4, fill=green!10] (check) {检验 + 灵敏度};
    \draw[arrow] (data) -- (q1);
    \draw[arrow] (q1) -- node[right]{指标/参数} (q2);
    \draw[arrow] (q2) -- node[right]{预测值} (q3);
    \draw[arrow] (q3) -- node[right]{分组结果} (q4);
    \draw[arrow] (q4) -- (check);
  \end{tikzpicture}
  \caption{本文整体求解框架}
  \label{fig:overall_route}
\end{figure}

% 跑题后必须删下面这一行 (check 15 占位符自检)
\textbf{【这里粘贴从第 1 段到第 5 段拼起来的不超过 2 页内容, 删掉这一行】}

%% ==================== 模型假设 模板 (v1.5.6, BZD 数模社 2026 规范) ====================
%%
%% 借鉴 BZD 数模社 第三章 模板: 引导语 + 4-8 条假设 + 单条 "假设 + 说明"
%%
%% BZD 9 类假设:
%%  1. 数据真实性与有效性 (数据集完整且不存在错误, 实际是 "经清洗后保留的数据满足建模需要")
%%  2. 小概率事件排除 (战争/重大灾害等)
%%  3. 主要因素保留、次要因素忽略
%%  4. 系统平稳性 (环境/参数/状态在范围内稳定)
%%  5. 对象理想化 (质点/圆柱/均匀混合)
%%  6. 参数形式与概率分布 (服从某分布)
%%  7. 变量关系 (线性/可加/其他)
%%  8. 行为与决策 (决策者准则)
%%  9. 空间与时间 (同一周期口径一致)
%% 10. 独立性或无相互干扰
%%
%% 写法: 每条 "假设 + 说明 (依据/简化/作用)".
%% 关键假设建议 1-3 句话说明理由, 简单题目给的假设可简短.
%% 数量: 4-8 条最佳, 不硬性.
%%
%% 跑题用户必须:
%% 1. 删掉所有 % 注释 (整段 L1-L23)
%% 2. 替换【】占位符为自己的实际内容
%% 3. 删 下面的 \AssumptionSlot{...} 行 (check 15 抓这个 inline 【】)

\section{模型假设}

为突出研究对象的主要特征并降低次要因素对模型求解的干扰, 在不改变问题本质的前提下, 结合题目条件、数据特征及后续模型需要, 作出如下假设:

\begin{enumerate}[label=(\arabic*)]

\item \textbf{【假设 1 内容, 1 句。\textbf{粗体}突出关键术语】}。

\textbf{说明:} 【假设依据 + 简化对象 + 建模作用, 1-3 句。例如: "数据均来自题目附件, 经缺失值和异常值处理后保留的数据能够反映研究对象的基本特征。该假设是数据预处理和模型输入的合理性前提, 不影响后续建模结论。"】

\item \textbf{【假设 2 内容, 1 句】}。

\textbf{说明:} 【依据 + 简化对象 + 建模作用, 1-3 句】

\item \textbf{【假设 3 内容, 1 句】}。

\textbf{说明:} 【依据 + 简化对象 + 建模作用, 1-3 句】

\item \textbf{【假设 4 内容, 1 句】}。

\textbf{说明:} 【依据 + 简化对象 + 建模作用, 1-3 句】

\item \textbf{【假设 5 内容, 1 句】}。

\textbf{说明:} 【依据 + 简化对象 + 建模作用, 1-3 句】

\item \textbf{【假设 6 内容, 1 句 (可选)】}。

\textbf{说明:} 【依据 + 简化对象 + 建模作用, 1-3 句】

\end{enumerate}

% 不推荐开头: "为了方便模型建立, 使模型更加完备, 预测结果更加合理, 作出如下假设。" -- 空泛.
% 不可取: 假设模型必然正确 / 假设数据绝对没有缺失 / 假设后文不用的假设 / 重复近义句.

% 跑题后必须删下面这一行 (check 15 占位符自检)
\textbf{【这里粘贴从假设 1 到假设 6 (4-8 条) 拼起来的不超过 1 页内容, 删掉这一行】}

%% ==================== 符号说明 模板 (v1.5.6, BZD 数模社 2026 规范) ====================
%%
%% 借鉴 BZD 数模社 第四章 模板: 引导语 + 三线表 (符号/含义/单位)
%%
%% BZD 原则:
%% - 三线表只保留顶线、表头分隔线、底线 (无左右竖线, 无逐行横线)
%% - "符号"列较窄居中, "含义"列最宽左对齐, "单位"列较窄居中
%% - 无量纲统一标 "—" 或 "无量纲" (全文选一种)
%% - 数量: 通常 8-15 个核心全局符号, 半页至一页
%% - min/max/sum/ln/exp 等常见运算符不入表 (除非论文反复用)
%% - 同一变量只用一个符号, 同一符号只表示一个含义
%%
%% 跑题用户必须:
%% 1. 删掉所有 % 注释 (整段 L1-L21)
%% 2. 替换【】占位符为自己的实际符号/含义/单位
%% 3. 删 下面的 \SymbolSlot{...} 行 (check 15 抓这个 inline 【】)

\section{符号说明}

为便于后续模型的建立、求解与结果分析, 现将全文反复使用的主要符号及其含义汇总如表 \ref{tab:symbols} 所示。各模型中仅局部使用的符号将在首次出现时另行说明。

\begin{table}[htbp]
  \centering
  \caption{主要符号说明}
  \label{tab:symbols}
  \begin{tabular}{p{0.15\textwidth} p{0.55\textwidth} p{0.18\textwidth}}
    \toprule
    \textbf{符号} & \textbf{含义} & \textbf{单位} \\
    \midrule
    【$x_i$】 & 【第 $i$ 个样本的观测值】 & 【—】 \\
    【$w_j$】 & 【第 $j$ 个指标的权重】 & 【—】 \\
    【$\mu$】 & 【样本均值】 & 【按变量确定】 \\
    【$\sigma$】 & 【样本标准差】 & 【按变量确定】 \\
    【$n$】 & 【样本量】 & 【—】 \\
    【$m$】 & 【指标数量】 & 【—】 \\
    【$t$】 & 【时间变量】 & 【年/月/日】 \\
    【$R^2$】 & 【决定系数, 衡量模型拟合优度】 & 【—】 \\
    【$\text{MAE}$】 & 【平均绝对误差】 & 【按变量确定】 \\
    【$\text{RMSE}$】 & 【均方根误差】 & 【按变量确定】 \\
    【$\text{AUC}$】 & 【ROC 曲线下面积】 & 【—】 \\
    【$\alpha$】 & 【显著性水平】 & 【—】 \\
    【$P$】 & 【概率】 & 【—】 \\
    \bottomrule
  \end{tabular}
\end{table}

\noindent\textbf{注:} 表中未列出的局部符号将在正文首次出现时说明。

% 不收录: 临时变量 (只在公式附近解释) / 局部符号 (该小节首次出现时说明) / 常见运算符.
% 必查: 一个符号只对应一个含义 / 同一变量只用一个符号 / 大小写和上下标一致 / 单位和正文一致.

% 跑题后必须删下面这一行 (check 15 占位符自检)
\textbf{【这里粘贴表格内容 (10-15 个核心符号), 删掉这一行】}

%% ==================== 问题 X 建模思路与数据准备 模板 (v1.5.6, BZD 数模社 2026 规范) ====================
%%
%% 借鉴 BZD 数模社 第五章 模板: 5.X.1 建模思路 + 数据处理
%%
%% BZD 结构 (5.X.1):
%%   - 建模思路: 一段话说明任务 + 数据特征 + 数学类型 + 模型选择理由
%%   - 数据处理: 数据来源 + 处理方法 + 处理后保留样本 + 特征构造
%%   - 算法选择: 模型/算法名称 + 选择依据
%%
%% 跑题用户必须:
%% 1. 删掉所有 % 注释
%% 2. 替换【】占位符为自己的实际内容
%% 3. 删 下面的 \ModelingSlot{...} 行 (check 15 抓这个 inline 【】)

\subsubsection{问题 X 建模思路与数据准备}

\textbf{（1）建模思路} \quad 问题 X 要求【问题目标, 1 句。例如: 分析影响 Y 染色体浓度稳定性的关键因素】。根据题目条件和数据特征, 该问题在数学上属于【数学类型。例如: 影响因素识别 + 多元回归分析】。考虑到【数据特点, 1 句。例如: 特征维度高 (10 个变量) 且变量间可能存在共线性】, 选择【主模型, 如 多元线性回归 (MLR) + 随机森林回归 (RFR) 组合】进行求解, 并针对【本题特点】对模型的【目标函数/约束/参数/算法】进行调整。

\textbf{（2）数据处理} \quad 本问题使用【数据来源, 1 句。例如: 题目附件 Sheet1 男胎 1082 例】。首先进行【处理方法, 1 句。例如: 缺失值检测 + 异常值剔除 (3$\sigma$ 准则) + 标准化 (z-score)】, 原因是【处理理由, 1 句。例如: 保证数据质量并消除量纲影响】。处理后保留【样本数量】个有效样本, 并构造【特征名称, 1 句】作为模型输入。

\textbf{（3）算法选择与依据} \quad 为保证模型可解释性, 同时具备非线性捕捉能力, 本文采用【主模型 + 备选模型】组合方案:
\begin{itemize}
  \item \textbf{主模型}:【如: 多元线性回归 (MLR), 数学类型为线性回归, 用于识别主要线性影响因素】
  \item \textbf{备选模型}:【如: 随机森林回归 (RFR), 集成 100 棵决策树, 用于捕捉非线性关系并提供特征重要性排序】
  \item \textbf{选择依据}:【如: MLR 可解释性强便于系数分析, RFR 拟合优度通常更高 (R²) 且提供特征重要性, 两者对比可避免单一模型的局限】
\end{itemize}

\textbf{（4）模型评价指标} \quad 本文采用以下指标评价模型表现:
\begin{itemize}
  \item $R^2$ (决定系数): 衡量模型解释因变量变异的能力, 越接近 1 越好
  \item $\text{RMSE}$ (均方根误差): 衡量预测值与真实值的偏差, 越小越好
  \item $\text{MAE}$ (平均绝对误差): 衡量预测值与真实值的绝对偏差, 越小越好
\end{itemize}

% 跑题后必须删下面这一行 (check 15 占位符自检)
\textbf{【这里粘贴建模思路 + 数据处理 + 算法选择 拼起来的不超过 1 页内容, 删掉这一行】}

%% ==================== 问题 X 建模与求解 模板 (v1.5.6, BZD 数模社 2026 规范) ====================
%%
%% 借鉴 BZD 数模社 第五章 模板: 5.X.2 建模思路 + 5.X.3 模型建立 + 5.X.4 求解 + 5.X.5 结果分析
%%
%% BZD 5.X.2-5.X.5 结构:
%%   - 模型建立: 定义变量 → 数学表达 → 目标函数/约束 → 参数说明
%%   - 模型求解: 算法步骤 → 关键参数 → 求解结果
%%   - 结果分析: 数值 + 单位 + 现实意义 + 异常解释
%%
%% 跑题用户必须:
%% 1. 删掉所有 % 注释
%% 2. 替换【】占位符为自己的实际内容
%% 3. 删 下面的 \SolvingSlot{...} 行 (check 15 抓这个 inline 【】)

\subsubsection{问题 X 模型的建立}

\textbf{（1）变量与参数定义} \quad 设【决策变量/状态变量/输入变量/输出变量】为:
\begin{itemize}
  \item $x_i$ : 【第 $i$ 个样本的【类型】观测值】
  \item $y_j$ : 【第 $j$ 个【目标类型】】
  \item $\beta_k$ : 【第 $k$ 个回归系数, 表示【类型】对【目标类型】的影响程度】
  \item $\varepsilon_i$ : 【第 $i$ 个随机扰动项, 满足 $\mathbb{E}(\varepsilon_i)=0$】
\end{itemize}

\textbf{（2）核心数学表达} \quad 根据【题目条件 + 问题类型】, 建立如下【模型类型】:
\begin{equation}
  y = f(x; \theta) + \varepsilon
  \label{eq:model}
\end{equation}

其中, $f(\cdot)$ 为【模型函数, 如 多元线性回归的线性组合 / 随机森林的非线性映射】, $\theta = (\beta_0, \beta_1, \ldots, \beta_p)$ 为待估参数向量, $\varepsilon$ 为随机扰动项。

对于【优化/预测/分类】任务, 综合考虑【目标要求】和【资源/时间/容量等限制】, 建立如下优化模型:
\begin{align}
  \min_{\theta} \quad & L(\theta) = \sum_{i=1}^{n} (y_i - f(x_i; \theta))^2 \\
  \text{s.t.} \quad & g_j(\theta) \leq 0, \quad j=1,\ldots,m \\
  & h_k(\theta) = 0, \quad k=1,\ldots,p \label{eq:constraints}
\end{align}

其中, $L(\theta)$ 为损失函数 (如平方损失/对数损失/铰链损失), $g_j$ 为不等式约束 (如非负、容量、阈值), $h_k$ 为等式约束 (如归一化、平衡条件)。

\textbf{（3）参数来源说明} \quad 模型参数来源于:
\begin{itemize}
  \item 题目附件直接给出 (如 BMI、孕周、染色体浓度阈值)
  \item 由数据回归或拟合获得 (如线性回归系数 $\beta$, 通过最小二乘求解)
  \item 通过优化算法标定 (如随机森林超参数, 通过网格搜索 + 交叉验证)
  \item 参考文献或领域经验 (如风险权重 $w_1:w_2 = 2:1$ 基于临床损失矩阵)
\end{itemize}

\textbf{（4）约束条件} \quad 每个约束的物理含义:
\begin{itemize}
  \item 【如: $g_1: \beta_k \geq 0$ 表示【主要因素对【目标】应有正向影响】
  \item 【如: $g_2: \sum_{i=1}^{n} w_i = 1$ 表示【权重归一化, 满足概率公理】
  \item 【如: $h_1: y_{pred} \in [0,1]$ 表示【预测概率在合理区间】
\end{itemize}

\subsubsection{问题 X 模型的求解}

\textbf{（1）算法选择} \quad 本文选择【算法名称, 如 序列最小二乘规划 (SLSQP) / 差分进化 (DE) / 随机森林集成 / 蒙特卡洛模拟】, 原因是【适用理由, 1 句】。

\textbf{（2）求解步骤}
\begin{enumerate}
  \item \textbf{数据准备}:【数据来源 + 预处理】, 形成输入矩阵 $X \in \mathbb{R}^{n \times p}$ 和目标向量 $y \in \mathbb{R}^n$
  \item \textbf{参数初始化}:【如 $\beta^{(0)} = (0,\ldots,0)$, 种群规模 N=100, 最大迭代 $T_{max}=500$】
  \item \textbf{迭代求解}:【如使用 scipy.optimize.minimize / sklearn.ensemble.RandomForestRegressor】
  \item \textbf{收敛判断}:【如 $\Delta L < 10^{-6}$ 或 $T = T_{max}$】
  \item \textbf{多次独立运行}:【为消除随机性, 重复运行 10 次取平均值, 随机种子固定】
\end{enumerate}

\textbf{（3）软件与版本} \quad Python 3.12 + numpy 1.26 + pandas 2.0 + scikit-learn 1.4 + scipy 1.11 + matplotlib 3.8。完整可运行代码见支撑材料 \texttt{求解/问题X/问题X\_xxx.py}。

\subsubsection{问题 X 求解结果与分析}

\textbf{（1）模型拟合结果} \quad 模型对训练集的拟合效果: $R^2 = 0.XXXX$, $\text{RMSE} = 0.XXXX$, $\text{MAE} = 0.XXXX$。对测试集的泛化能力: $R^2 = 0.XXXX$, $\text{RMSE} = 0.XXXX$, $\text{MAE} = 0.XXXX$。【如 随机森林回归明显优于多元线性回归, 表明变量间存在非线性关系】

\textbf{（2）关键参数估计} \quad 【主要因素】的回归系数 $\beta_1 = 0.XXXX$ ($p < 0.001$), 表明【主要因素】对【目标】有【正向/负向】显著影响。【次要因素】的系数 $\beta_2 = 0.XXXX$ ($p > 0.05$), 不显著, 在后续模型中可剔除。

\textbf{（3）结果与现实一致性} \quad 本文得到的【主要结论, 如: 孕周每增加 1 周, Y 染色体浓度平均增加 0.0074】与临床观察一致。BMI 的负向影响 (系数 -0.0015) 符合【领域知识, 如: 母体 BMI 越高, 胎儿游离 DNA 浓度越低】的预期。

\textbf{（4）可视化} \quad 如图所示, 【主要可视化内容, 如: Y 染色体浓度 vs 孕周散点图 + 回归拟合曲线 / 特征重要性柱状图 / 残差分布直方图】。

% 跑题时由用户在此处补 figure 环境:
% \begin{figure}[htbp]
%   \centering
%   \includegraphics[width=0.9\textwidth]{figures/q1_特征重要性.png}
%   \caption{问题 X 特征重要性排序}
%   \label{fig:q1_results}
% \end{figure}

% 跑题后必须删下面这一行 (check 15 占位符自检)
\textbf{【这里粘贴问题 X 完整的 5.X.2-5.X.5 内容, 包括模型建立 + 求解 + 结果分析, 不超过 6-7 页, 删掉这一行】}

%% ==================== 模型检验 模板 (v1.5.6, BZD 数模社 2026 规范) ====================
%%
%% 借鉴 BZD 数模社 模型检验模板: 误差分析 + 交叉验证 + 灵敏度 + 稳健性
%%
%% BZD 模型检验 4 大类:
%%  1. 误差分析: MAE/RMSE/MAPE/R² + 残差分布
%%  2. 交叉验证: K 折 (通常 K=5 或 K=10) 验证泛化能力
%%  3. 灵敏度分析: 关键参数 ±10-20% 变化对结果影响
%%  4. 稳健性: 改变输入数据/异常值处理/参数初值, 看结论是否改变
%%
%% 跑题用户必须:
%% 1. 删掉所有 % 注释
%% 2. 替换【】占位符为自己的实际内容
%% 3. 删 下面的 \CheckSlot{...} 行 (check 15 抓这个 inline 【】)

\section{模型检验}

\subsection{误差分析}

为评估各模型的预测精度, 本文采用 MAE (平均绝对误差)、RMSE (均方根误差)、MAPE (平均绝对百分比误差) 和 $R^2$ (决定系数) 四个指标, 在测试集上对各模型进行评价, 结果如表 \ref{tab:error_metrics} 所示。

\begin{table}[htbp]
  \centering
  \caption{各模型在测试集上的误差指标}
  \label{tab:error_metrics}
  \begin{tabular}{lcccc}
    \toprule
    \textbf{模型} & \textbf{MAE} & \textbf{RMSE} & \textbf{MAPE (\%)} & \textbf{$R^2$} \\
    \midrule
    【主模型 (问题一)】 & 【0.0XXX】 & 【0.0XXX】 & 【X.XX】 & 【0.XXXX】 \\
    【主模型 (问题二)】 & 【0.0XXX】 & 【0.0XXX】 & 【X.XX】 & 【0.XXXX】 \\
    【主模型 (问题三)】 & 【0.0XXX】 & 【0.0XXX】 & 【X.XX】 & 【0.XXXX】 \\
    【主模型 (问题四)】 & 【0.0XXX】 & 【0.0XXX】 & 【X.XX】 & 【0.XXXX】 \\
    \bottomrule
  \end{tabular}
\end{table}

从表 \ref{tab:error_metrics} 可以看出, 【主结论, 如: 本文建立的【模型名】在【数据范围】内具有较好的预测精度, 各项误差指标均处于可接受水平, $R^2$ 均在 0.XX 以上, 表明模型能够解释【目标变量】约 XX\% 的变异】。

\subsection{交叉验证}

为验证模型的泛化能力, 本文采用 5 折交叉验证 (K=5), 将【总样本量 N】个样本随机划分为 5 个子集, 每次以 4 个子集训练、1 个子集测试, 循环 5 次, 取平均结果。交叉验证结果如表 \ref{tab:cv} 所示。

\begin{table}[htbp]
  \centering
  \caption{5 折交叉验证结果}
  \label{tab:cv}
  \begin{tabular}{cccccc}
    \toprule
    \textbf{折次} & \textbf{1} & \textbf{2} & \textbf{3} & \textbf{4} & \textbf{5} \\
    \midrule
    $R^2$ & 【0.XXXX】 & 【0.XXXX】 & 【0.XXXX】 & 【0.XXXX】 & 【0.XXXX】 \\
    \bottomrule
  \end{tabular}
\end{table}

5 折交叉验证的 $R^2$ 标准差为 【0.0XX】, 表明模型在【数据规模】下具有良好的稳定性和泛化能力, 不存在明显过拟合。

\subsection{灵敏度分析}

为检验模型对关键参数的敏感程度, 本文选取【主要参数, 2-3 个】, 在【基准值 $\pm 20\%$】范围内逐步变化, 观察模型输出的变化情况, 结果如图 \ref{fig:sensitivity} 所示。

\begin{figure}[htbp]
  \centering
  \includegraphics[width=0.9\textwidth]{figures/sensitivity_关键参数.png}
  \caption{关键参数灵敏度分析 (基准值 ± 20\%)}
  \label{fig:sensitivity}
\end{figure}

灵敏度分析表明:
\begin{itemize}
  \item 【参数 1】在【区间】内变化时, 模型结果保持稳定, 相对灵敏度系数 $|S| < 0.1$, 表明模型对该参数不敏感
  \item 【参数 2】在【区间】内变化时, 模型结果变化明显, 相对灵敏度系数 $|S| > 1.0$, 表明模型对该参数较为敏感, 在实际应用中应【提高该参数的估计精度 / 增加额外验证 / 设置保守阈值】
\end{itemize}

\subsection{稳健性检验}

为检验模型对数据扰动的稳健性, 本文采用以下方式:
\begin{itemize}
  \item \textbf{异常值处理方法变化}: 分别使用 3$\sigma$ 准则和 IQR 准则 (四分位距) 剔除异常值, 重新训练模型, 主要结论保持不变
  \item \textbf{样本规模变化}: 随机抽取 70\% 和 50\% 的样本重新训练, 关键参数的相对变化均小于 10\%, 表明模型对样本规模不敏感
  \item \textbf{参数初值变化}: 对随机算法 (如随机森林、差分进化) 使用不同随机种子重复运行 10 次, 关键结果的标准差均小于 5\%, 表明算法具有较好的稳定性
\end{itemize}

% 跑题后必须删下面这一行 (check 15 占位符自检)
\textbf{【这里粘贴完整的 4 类检验结果, 不超过 2-3 页, 删掉这一行】}

%% ==================== 模型评价 模板 (v1.5.6, BZD 数模社 2026 规范) ====================
%%
%% 借鉴 BZD 数模社 第六章 模板: 总结 + 优点 + 不足
%% 注意: 4 子节结构 (BZD 常见形式 A) 或 3 子节 (形式 B)
%%
%% BZD 原则:
%% - 优点: "具体做法 + 形成的优势 + 证据或效果" (避免 "精度高、创新性强" 套话)
%% - 不足: "不足来源 + 可能影响 + 影响范围" (避免足以否定全文的致命问题)
%% - 改进: 必须逐项对应不足
%% - 推广: 说明对象 + 变量替换 + 参数标定 + 约束调整 + 所需数据
%%
%% 跑题用户必须:
%% 1. 删掉所有 % 注释
%% 2. 替换【】占位符为自己的实际内容
%% 3. 删 下面的 \EvalSlot{...} 行 (check 15 抓这个 inline 【】)

\section{模型评价、改进与推广}

\subsection{模型总结}

本文针对【研究问题】, 分别建立【模型 1】、【模型 2】、【模型 3】和【模型 4】。各模型依次完成【问题一目标】、【问题二目标】、【问题三目标】和【问题四目标】, 并通过【检验方法: 5 折交叉验证 + 灵敏度分析 + 稳健性检验】验证结果可靠性。

问题一通过【方法, 如 多元线性回归 + 随机森林回归】识别【主要因素】; 问题二在问题一的基础上建立【如 ROC 曲线 + 贝叶斯决策】, 推荐【最佳时点】; 问题三在问题二的基础上按【分组标准】分层, 揭示【分组差异】; 问题四在问题三的基础上建立【如 随机森林 + Isolation Forest】, 形成【决策规则 + 比例】。

经检验, 各模型在【数据范围】内具有较好的【准确性 ($R^2 > 0.XX$)、稳定性 (CV 标准差 < 0.XX) 和泛化能力 (测试集误差 < 训练集)】, 可应用于【具体场景或推广对象】。

\subsection{模型的优点}

（1）\textbf{模型具有较强针对性}。针对【题目特点】, 在【经典模型】基础上增加【新变量/新约束/新参数】, 提高了模型与实际问题的匹配程度。【如: 在标准 Logistic 回归基础上引入 BMI $\times$ 孕周交互项, 揭示了 BMI 对 NIPT 检测时点的非线性影响】

（2）\textbf{数据处理较为完整}。针对【数据问题】采用【处理方法: 缺失值 + 异常值 + 标准化】, 降低了数据质量对后续模型的影响。【如: 经清洗后保留 N 个有效样本, 处理后的数据通过了 K-S 正态性检验和方差齐性检验】

（3）\textbf{问题之间联动清晰}。将问题一的【输出: 主要特征/权重/参数】作为问题二的【输入】, 进一步进入问题三的【目标函数或约束】, 最后由问题四形成【综合决策】, 实现了数据流、参数流和结论流的完整串联。【如: 特征重要性 (0.42) 作为问题二的关键特征, 预测值作为问题三的输入参数, 决策概率作为问题四的判据】

（4）\textbf{模型结果可靠}。经【检验方法: 5 折交叉验证 + 灵敏度分析 + 稳健性检验】验证, 【量化指标, 如: $R^2 = 0.XX$, $\text{AUC} = 0.XX$】达到【具体数值】, 主要结论与【现实规律或临床观察】一致。

（5）\textbf{具有推广价值}。调整【参数或约束】后, 可用于【相似对象或场景】, 详见 7.4 节。

\subsection{模型的不足}

（1）受【数据来源: 单一医院 / 单一地区 / 时间跨度短】限制, 未能充分考虑【其他因素】, 可能对【具体结果】产生一定影响。【如: 本文数据来自单家三甲医院, 缺乏多中心验证, 推广到基层医院或不同地区时可能存在偏差】

（2）为保证模型可解, 对【现实过程】进行了【具体简化】, 使模型在【特定情景】下的适用性受到限制。【如: 假设各 BMI 组内部完全独立, 但实际可能存在组间空间相关或遗传背景相似】

（3）模型结果对参数【参数名】较为敏感, 该参数估计误差可能影响最终结论。【如: 风险权重 $w_1:w_2 = 2:1$ 基于临床经验, 不同权重设置下最优方案可能略有变化, 但整体结论 (最佳时点) 保持稳定】

（4）【算法名称】存在一定随机性, 虽然多次独立运行结果较稳定, 但不能从理论上保证全局最优。【如: 差分进化算法可能陷入局部最优, 可通过增加种群规模或迭代次数改善】

\subsection{模型的改进与推广}

\textbf{（1）模型改进方向} \quad 针对上述不足, 本文提出以下改进方向:
\begin{itemize}
  \item 进一步收集【数据类型, 如 多中心数据 / 更长时间序列】, 扩大时间和空间范围, 提高参数估计可靠性
  \item 增加【新变量/新约束】, 进一步描述当前未充分考虑的现实因素
  \item 引入【动态模型 / 鲁棒优化 / 改进算法】, 降低不确定性对结果的影响
  \item 使用【更精细的随机性控制或理论保证】, 解决算法的随机性问题
\end{itemize}

\textbf{（2）模型推广} \quad 本文模型的核心结构是根据【输入变量】建立【目标变量】与【约束】之间的关系, 因此可以推广至【其他对象, 如 类似疾病筛查 / 类似临床决策问题】。推广时需要:
\begin{itemize}
  \item 将【原变量, 如 Y 染色体浓度】替换为【新变量, 如 其他疾病相关游离 DNA 浓度】
  \item 重新标定【关键参数, 如 BMI 阈值、风险权重】
  \item 增加【新约束, 如 不同疾病的阈值和决策规则】
  \item 补充【新数据, 如 目标疾病的检测数据】
\end{itemize}

推广后的模型可用于解决【类似任务, 如 其他染色体非整倍体筛查 / 其他无创产前检测】, 具有较好的实际应用价值。

% 跑题后必须删下面这一行 (check 15 占位符自检)
\textbf{【这里粘贴完整的 4 子节内容, 不超过 1-2 页, 删掉这一行】}

%% ==================== 进一步工作方向 模板 (v1.5.6, BZD 数模社 2026 规范) ====================
%%
%% 借鉴 BZD 数模社 第七章 模型评价、改进与推广 的 "改进" 和 "推广" 子节展开
%% 注意: 7.模型评价.tex 已经包含简明总结 + 优点 + 不足 + 改进与推广 4 子节.
%%       本章 8 是 "进一步工作方向", 比 7 更具体的待做工作, 可作为赛后延伸.
%%
%% BZD 改进方向 (从 7.3 拓展):
%% - 数据: 扩大样本/多中心/外部验证
%% - 模型: 增加变量/非参数/动态/多目标/随机规划
%% - 算法: 改进初始化/混合算法/并行计算/降维
%%
%% BZD 推广方向:
%% - 对象: 推广到相似问题
%% - 区域: 从单点扩展到多地区
%% - 方法: 从单目标到多目标
%%
%% 跑题用户必须:
%% 1. 删掉所有 % 注释
%% 2. 替换【】占位符为自己的实际内容
%% 3. 删 下面的 \FutureSlot{...} 行 (check 15 抓这个 inline 【】)

\section{进一步工作方向}

\subsection{数据层面的改进}

\begin{enumerate}
  \item \textbf{扩大样本规模和地域覆盖}。本文数据来自【来源, 如 单一医院 N 例】, 后续可联合【多中心 / 跨地区】医院, 收集更大规模的数据, 验证模型在不同人群中的普适性, 同时进行【外部验证 (External Validation)】。

  \item \textbf{增加数据类型和特征维度}。除【现有变量】外, 可补充【新特征, 如 基因型、孕期并发症、胎儿性别、孕产次】, 进一步提升模型预测能力。

  \item \textbf{建立长期随访数据}。建立【长期队列, 如 5-10 年】, 跟踪【研究对象】的长期结局, 验证模型的长期预测准确性和实际应用价值。
\end{enumerate}

\subsection{模型层面的改进}

\begin{enumerate}
  \item \textbf{引入非参数和深度学习方法}。对于【非线性关系复杂】的数据, 可尝试【如 神经网络、Transformer、贝叶斯神经网络】等深度学习方法, 进一步提升模型拟合能力。

  \item \textbf{构建多目标优化模型}。在【如 漏检与过度诊断的权衡】等多目标场景下, 可建立多目标优化模型, 使用【NSGA-II / MOEA/D】等算法, 提供帕累托前沿供决策者选择。

  \item \textbf{引入不确定性量化}。使用【如 贝叶斯神经网络、深度集成、MC Dropout】等方法, 给出预测结果的置信区间, 为临床决策提供更全面的信息。

  \item \textbf{动态建模与时序分析}。考虑【如 孕周、BMI 随时间的动态变化】, 建立【动态模型 / 状态空间模型 / 循环神经网络】, 刻画研究对象的动态演化规律。
\end{enumerate}

\subsection{算法层面的改进}

\begin{enumerate}
  \item \textbf{改进优化算法的初始化和参数}。对【差分进化 / 粒子群 / 遗传算法】等智能优化算法, 使用【拉丁超立方 / 对立学习 / Tent 混沌】等改进初始化策略, 提升全局搜索能力。

  \item \textbf{采用混合算法}。将【如 差分进化 + 局部搜索 / 模拟退火】结合, 兼顾全局探索和局部开发, 提升解的质量和收敛速度。

  \item \textbf{并行化与高性能计算}。对【如 蒙特卡洛模拟 / 交叉验证 / 灵敏度分析】等计算密集型任务, 使用【MPI / GPU / 多线程】并行化, 显著缩短计算时间。
\end{enumerate}

\subsection{应用层面的推广}

\begin{enumerate}
  \item \textbf{推广到其他染色体非整倍体疾病}。本文以【NIPT 13/18/21】为例建立模型, 该方法可推广到【如 性染色体非整倍体 (Turner, Klinefelter)、微缺失综合征 (DiGeorge, 1p36)】等疾病的产前筛查。

  \item \textbf{推广到其他临床决策问题}。本文的【如 风险决策 + 阈值优化】方法可推广到【肿瘤标志物诊断、心血管疾病风险预测、传染病筛查】等其他临床决策问题。

  \item \textbf{开发临床决策支持系统 (CDSS)}。将本文模型封装为【可解释的临床决策支持工具, 如 Web 应用 / 移动端 App】, 辅助基层医生进行【产前咨询 + 最佳检测时点推荐】, 推动模型的临床转化。
\end{enumerate}

% 跑题后必须删下面这一行 (check 15 占位符自检)
\textbf{【这里粘贴从数据/模型/算法/应用 4 个层面的进一步工作, 不超过 1 页, 删掉这一行】}


\newpage
%% 2026 规范：AI 工具使用声明必须在参考文献之前
%% ==================== AI 工具使用声明 模板 (v1.5.6, BZD 数模社 2026 规范) ====================
%%
%% 借鉴 BZD 数模社 论文末尾模板: 二者择一 (未使用 AI / 使用 AI)
%% 必须在参考文献之前, AI 工具不列入参考文献, 详细情况见支撑材料《AI工具使用详情.pdf》
%%
%% 跑题用户必须:
%% 1. 删掉所有 % 注释
%% 2. 二者择一: 删除其中一段, 保留另一段
%% 3. 替换【】占位符为自己的实际内容
%% 4. 删 下面的 \AISlot{...} 行 (check 15 抓这个 inline 【】)

\subsection*{AI 工具使用声明}

% ====== 二者择一: 跑题用户必须删掉其中一段, 保留另一段 ======

% ===== 选项 (1) 未使用 AI 工具: 保留下面这行, 删掉 "选项 (2)" =====
% 本参赛队在竞赛过程中未使用任何 AI 工具。

% ===== 选项 (2) 使用 AI 工具: 删掉 "选项 (1)" 上面这行, 保留下面这段, 并填实际用途 =====
%% ⚠️ §9 论文内 = 1 句话简短声明, 只填 3-5 项 **简要** 用途 (e.g. 读题 / 代码 / 公式 / 文档 / 可视化)
%%    **不要在这里展开 4 节详写**, 详细 5+ 项看 支撑材料/AI工具使用详情.pdf (4 节)
本参赛队在竞赛过程中使用了 AI 工具, 主要用于【简要填 3-5 项实际用途, e.g. 读题 / 代码 (LP + CBC) / 公式推导 / 文档润色 / 数据可视化】, 详细使用情况见支撑材料《AI 工具使用详情.pdf》。

% 跑题后必须删下面这一行 (check 15 占位符自检)
\AISlot{【这里填实际使用情况, 跑题后删这一行】}

% AI 工具不能列入参考文献! ChatGPT / DeepSeek / Claude / Copilot 等条目禁止.
% AI 推荐的文献只列原始论文 (经人工核验), 不引用 AI.
% 声明应与支撑材料《AI工具使用详情.pdf》内容一致.


%% ==================== 参考文献 模板 (v1.5.6, BZD 数模社 2026 规范) ====================
%%
%% 借鉴 BZD 数模社 参考文献模板: 5-15 条, GB/T 7714 格式
%%
%% BZD 原则:
%% - 顺序编码制, 按正文第一次引用的顺序编号为 [1] [2] [3]...
%% - 真实可核验, AI 推荐的文献必须由人工核对作者/题名/期刊/年份
%% - AI 工具 (ChatGPT/DeepSeek/Copilot) 禁止列入参考文献
%% - 5-15 条最佳, 优先经典 + 近 3-5 年 + 数据来源 + 标准
%% - 每条格式严格: 作者. 题名 [J/M/D/C/N/R/S/P/EB/OL/DB/OL/Z]. 出版地: 出版者, 年.
%%
%% 跑题用户必须:
%% 1. 删掉所有 % 注释
%% 2. 替换【】占位符为自己的实际参考文献
%% 3. 删 下面的 \RefSlot{...} 行 (check 15 抓这个 inline 【】)
%% 4. 按正文实际引用顺序重新编号

\section{参考文献}

% ====== 模板示例 (按 BZD 通用模板, 跑题后必须替换为真实文献) ======

\begin{enumerate}[label={[\arabic*]}, leftmargin=2em]
  \item \textbf{【作者 1】.} \textbf{【论文题名】} [J]. \textbf{【期刊名】}, \textbf{【年份】}, \textbf{【卷(期)】}: \textbf{【起止页码】}.
  \item \textbf{【作者 2】.} \textbf{【书名】} [M]. \textbf{【出版地】}: \textbf{【出版社】}, \textbf{【出版年】}: \textbf{【起止页码】}.
  \item \textbf{【作者 3】.} \textbf{【论文题名】} [D]. \textbf{【出版地】}: \textbf{【授予单位】}, \textbf{【出版年】}.
  \item \textbf{【发布机构】.} \textbf{【报告名称】} [R]. \textbf{【出版地】}: \textbf{【发布机构】}, \textbf{【出版年】}.
  \item \textbf{【发布机构】.} \textbf{【数据集名称】} [DB/OL]. \textbf{【网址】}, \textbf{【更新时间】} / \textbf{【访问日期】}.
  \item \textbf{【作者或机构】.} \textbf{【网页资源名称】} [EB/OL]. \textbf{【网址】}, \textbf{【发布日期】} / \textbf{【访问日期】}.
  \item \textbf{【标准号】.} \textbf{【标准名称】} [S]. \textbf{【出版地】}: \textbf{【出版者】}, \textbf{【出版年】}.
  \item \textbf{【颁布机构】.} \textbf{【文件名称】} [Z]. \textbf{【发布机构】}, \textbf{【发布日期】}.
\end{enumerate}

% 禁止: ChatGPT / DeepSeek / Claude / Copilot / 文心一言 / 通义千问 等 AI 工具条目.
% 禁止: AI 虚构作者 / 题名 / DOI / 期刊 / 页码 / 网址.
% 高风险: 正文有引用编号, 但文末没有对应条目 (一对一对应).
% 推荐: 经典模型原始文献 (Breiman 2001 Random Forest, Lo 1997 NIPT, Chiu 2008...) + 近 5 年相关 + 官方数据/标准.

% 跑题后必须删下面这一行 (check 15 占位符自检)
\RefSlot{【这里填 5-15 条真实可核验的参考文献, 按正文引用顺序编号, 删掉这一行】}


\newpage
%% 2026 规范第五/十一条：附录必须有「程序使用说明」+「支撑材料说明」两段固定句
%% ==================== 附录固定说明 模板 (v1.5.6, BZD 数模社 2026 规范) ====================
%%
%% 借鉴 BZD 数模社 附录模板: A 支撑材料清单 + B 软件环境 + C/D/E/F 源程序 + 补充
%%
%% 2026 规范 (第五/十一条 + 2026-08-26 官方答疑 Q3):
%% - 附录首页应是支撑材料清单 (file list)
%% - 附录 1: 程序使用说明
%% - 附录 2: 支撑材料说明
%% - 附录 3: 核心源代码 (推荐)
%% - 漏写 = 0 分 (评审硬性要求)
%%
%% 跑题用户必须:
%% 1. 删掉所有 % 注释
%% 2. 替换【】占位符为自己的实际文件清单和软件版本
%% 3. 删 下面的 \AppendixSlot{...} 行 (check 15 抓这个 inline 【】)

\section*{附录固定说明}

% ====== 附录 A: 支撑材料文件列表 (必填) ======
\subsection*{附录 A \quad 支撑材料文件列表}

\begin{table}[H]
  \centering
  \caption{支撑材料文件列表}
  \label{tab:appendix_files}
  \begin{tabular}{clll}
    \toprule
    \textbf{序号} & \textbf{文件名称} & \textbf{格式} & \textbf{主要内容} \\
    \midrule
    1 & \texttt{求解/问题1/问题1\_xxx.py} & PY & 问题一数据处理与模型求解 \\
    2 & \texttt{求解/问题1/图片/*.png} & PNG & 问题一输出图表 \\
    3 & \texttt{求解/问题1/结果/*.csv} & CSV & 问题一中间结果 \\
    4 & \texttt{求解/问题2/问题2\_yyy.py} & PY & 问题二预测与时点优化 \\
    5 & \texttt{求解/问题3/问题3\_zzz.py} & PY & 问题三分层优化 \\
    6 & \texttt{求解/问题4/问题4\_ddd.py} & PY & 问题四集成决策 \\
    7 & \texttt{求解/共享/constants.py} & PY & 跨问题物理常量 \\
    8 & \texttt{数据/原始数据.xlsx} & XLSX & 题目附件 (如有补充数据) \\
    9 & \texttt{AI工具使用详情.pdf} & PDF & AI 工具使用情况详细说明 \\
    \bottomrule
  \end{tabular}
\end{table}

% ====== 附录 B: 软件环境与求解工具 (必填) ======
\subsection*{附录 B \quad 软件环境与求解工具}

本文使用 \textbf{Python 3.12} 完成数据处理与模型求解, 主要程序库及版本:
\begin{itemize}
  \item numpy 1.26 (数值计算)
  \item pandas 2.0 (数据处理)
  \item scipy 1.11 (科学计算, 含 optimize 优化)
  \item scikit-learn 1.4 (机器学习, 含线性回归/随机森林/Isolation Forest)
  \item matplotlib 3.8 (可视化)
  \item PuLP 2.7 (线性规划, 如有优化问题)
\end{itemize}

使用 \textbf{Overleaf + xelatex (TeX Live 2024)} 完成论文排版, 使用 \textbf{TikZ} 绘制整体求解框架图。

% ====== 附录 C-F: 各问题源程序说明模板 (简化) ======
\subsection*{附录 C \quad 问题一完整源程序}

\begin{itemize}
  \item \textbf{程序名称}: 问题1\_xxx.py
  \item \textbf{编程语言及版本}: Python 3.12
  \item \textbf{对应问题}: 问题一
  \item \textbf{程序作用}: 【数据处理 / 模型求解 / 检验 / 绘图】
  \item \textbf{输入文件}: 【原始数据.xlsx 的 Sheet1】
  \item \textbf{输出文件}: 【问题1\_图片/*.png, 问题1\_结果/*.csv】
  \item \textbf{依赖环境}: numpy, pandas, scikit-learn
  \item \textbf{运行顺序}: \texttt{python 问题1/问题1\_xxx.py}
  \item \textbf{正文对应位置}: 第五章 5.1 节
\end{itemize}

完整源程序:
\begin{verbatim}
# 完整可运行的 Python 代码 (贴核心 30-50 行)
import numpy as np
import pandas as pd
from sklearn.ensemble import RandomForestRegressor

# 1. 加载数据
data = pd.read_excel('../../数据/原始数据.xlsx', sheet_name=0)

# 2. 数据预处理
X = data[['特征1', '特征2', '特征3']].values
y = data['目标变量'].values

# 3. 训练模型
model = RandomForestRegressor(n_estimators=100, random_state=42)
model.fit(X, y)

# 4. 预测与评估
y_pred = model.predict(X)
r2 = model.score(X, y)
print(f'R² = {r2:.4f}')
\end{verbatim}

\subsection*{附录 D \quad 问题二完整源程序}

\textbf{程序名称}: 问题2\_yyy.py

\textbf{编程语言及版本}: Python 3.12

\textbf{对应问题}: 问题二

\textbf{程序作用}: 【数据处理 / 模型求解 / 检验 / 绘图】

\textbf{输入文件}: 【原始数据.xlsx 的 Sheet1】+ 【问题一输出结果】

\textbf{输出文件}: 【问题2\_图片/*.png, 问题2\_结果/*.csv】

\textbf{依赖环境}: numpy, pandas, scikit-learn, scipy

\textbf{运行顺序}: \texttt{python 问题2/问题2\_yyy.py}

\textbf{正文对应位置}: 第五章 5.2 节

% 完整可运行代码见支撑材料对应文件

\subsection*{附录 E \quad 问题三完整源程序}

\textbf{程序名称}: 问题3\_zzz.py

\textbf{编程语言及版本}: Python 3.12

\textbf{对应问题}: 问题三

\textbf{程序作用}: 【数据处理 / 模型求解 / 检验 / 绘图】

\textbf{输入文件}: 【原始数据.xlsx 的 Sheet1】+ 【问题一/二输出结果】

\textbf{输出文件}: 【问题3\_图片/*.png, 问题3\_结果/*.csv】

\textbf{依赖环境}: numpy, pandas, scikit-learn, scipy

\textbf{运行顺序}: \texttt{python 问题3/问题3\_zzz.py}

\textbf{正文对应位置}: 第五章 5.3 节

% 完整可运行代码见支撑材料对应文件

\subsection*{附录 F \quad 问题四完整源程序}

\textbf{程序名称}: 问题4\_ddd.py

\textbf{编程语言及版本}: Python 3.12

\textbf{对应问题}: 问题四

\textbf{程序作用}: 【数据处理 / 模型求解 / 检验 / 绘图】

\textbf{输入文件}: 【原始数据.xlsx 的 Sheet1 或 Sheet2】+ 【问题一/二/三输出结果】

\textbf{输出文件}: 【问题4\_图片/*.png, 问题4\_结果/*.csv】

\textbf{依赖环境}: numpy, pandas, scikit-learn

\textbf{运行顺序}: \texttt{python 问题4/问题4\_ddd.py}

\textbf{正文对应位置}: 第五章 5.4 节

% 完整可运行代码见支撑材料对应文件

% 跑题后必须删下面这一行 (check 15 占位符自检)
\AppendixSlot{【这里填实际的文件清单 (附录 A 表格) + 软件环境 (附录 B) + 各问题源程序说明 (附录 C-F), 删掉这一行】}


%%%%%%%%%%%%%%%%%%%%%%%%%%%%%%%%%%%%%%%%%%%%%%%%%%%%%%%%%%%%%
% 通用附录模板 (v1.5.5 重写, 通用化) — 替代原 2024B 烟幕题 example
%
% 设计原则:
% 1. 整段都是 \TeX 注释 (以 % 开头), check 16 跳过 % 行 → dryrun 不报 RED
% 2. 留 1 行 INLINE 占位符 \AppendixSlot{...}, check 16 抓 \texttt{} 引用 → 跑题用户填完才 GREEN
% 3. 跑题用户复制本文件后, 删注释 + 填实际文件路径, 删 \AppendixSlot 行
%
% 2026 规范第五/十一条 + 2026-08-26 官方答疑 Q3:
% - 附录首页应是支撑材料清单 (file list)
% - 附录 1: 程序使用说明 (用了/没用)
% - 附录 2: 支撑材料说明 (提供了/没有)
% - 附录 3: 核心源代码 (推荐)
% - ⚠️ 漏写就是 0 分 (评审硬性要求)
%%%%%%%%%%%%%%%%%%%%%%%%%%%%%%%%%%%%%%%%%%%%%%%%%%%%%%%%%%%%%

\section*{附录 1：程序使用说明}

%%\subsection*{运行环境}
%%本参赛队使用 Python \texttt{3.12} + 主要依赖: numpy $\ge 1.24$, scipy $\ge 1.11$, pandas $\ge 2.0$, matplotlib $\ge 3.7$, openpyxl $\ge 3.1$。如果用了其他库 (如 scikit-learn, PuLP, PyMC), 在此列出。

%%\subsection*{目录结构}
%%\texttt{求解/} 目录下的主要文件:
%%\begin{itemize}
%%  \item \texttt{共享/constants.py} --- 跨问题物理常量
%%  \item \texttt{共享/params.py} --- 跨问题复用参数
%%  \item \texttt{问题1/问题1\_xxx.py} --- 问题 1 主代码
%%  \item \texttt{问题1/图片/*.png} --- 问题 1 输出图
%%  \item \texttt{问题1/结果/*.csv} --- 问题 1 输出表
%%  ... (按你实际问题列, 1-2 行/文件)
%%\end{itemize}

%%\subsection*{运行方法}
%%\begin{verbatim}
%%cd 求解
%%python 问题1/问题1_xxx.py
%%python 问题2/问题2_yyy.py
%%... (按你实际问题列)
%%\end{verbatim}

\section*{附录 2：支撑材料说明}

%%支撑材料包 \texttt{支撑材料.rar} 内容:
%%\begin{enumerate}
%%  \item \texttt{求解/}: 所有 Python 源代码
%%  \item \texttt{求解/问题*/图片/}: 所有图表 (PNG)
%%  \item \texttt{求解/问题*/结果/}: 所有结果 (CSV/XLSX)
%%  \item \texttt{AI工具使用详情.pdf}: AI 工具使用情况详细说明
%%  \item (如果用) \texttt{原始数据.xlsx}: 题面附件
%%\end{enumerate}
%%支撑材料大小: 约 X MB (远小于 20 MB 上限)。

\section*{附录 3：核心源代码}

%% 官方 Word 模板建议: 附录 3 放核心源代码 (贴关键文件 / 节选片段),
%% 配合"该代码使用 AI 工具: ..."标注。挑 1-2 个最关键函数贴入,
%% 不用整文件贴 (会超页)。

%%\subsection*{关键函数 1 (xxx.py 第 N 行)}
%%该部分代码使用 AI 工具: 【工具名, 版本, 日期】 / 团队手写
%%\begin{verbatim}
%%def your_key_function(...):
%%    """功能描述"""
%%    ... (贴核心逻辑, 30-50 行内)
%%\end{verbatim}

% ===== 跑题用户最终要填的实际内容 =====
% 把上面所有 % 注释替换为你的实际内容, 然后在下面填 1 行 inline 占位符
% 让 dryrun check 16 验证你的支撑材料文件清单是否都在

\textbf{【在这里列你实际的 .py / .png / .xlsx 文件路径, 跟支撑材料.zip 完全对齐。dryrun check 16 会逐个 Test-Path 验证。本行留【】= dryrun check 15 RED, 必须替换】}


\end{document}
